\documentclass[11pt]{article}
\usepackage[margin=0.85in]{geometry}
\usepackage{hyperref}
\usepackage{enumitem}
\setlist[itemize]{leftmargin=1.4em,itemsep=0.25em,topsep=0.35em}
\setlist[enumerate]{leftmargin=1.5em,itemsep=0.25em,topsep=0.35em}
\setlength{\parindent}{0pt}
\setlength{\parskip}{0.55em}
\begin{document}
\begin{center}
{\LARGE \textbf{Theory Report: Random matrix theory of integrability-to-chaos transition}}\\[0.4em]
{\large Ben Craps, Marine De Clerck, Oleg Evnin, Maxim Pavlov}\\[0.25em]
\href{https://arxiv.org/abs/2604.03669}{arXiv:2604.03669} \hspace{1em} Date: 2026-04-08
\end{center}

\section*{Paper information}
\textbf{Title.} Random matrix theory of integrability-to-chaos transition

\textbf{Authors.} Ben Craps, Marine De Clerck, Oleg Evnin, Maxim Pavlov

\textbf{Core question.} When an integrable quantum Hamiltonian is deformed by a chaotic perturbation, what statistical information determines the intermediate level-spacing law between the Poisson and Wigner--Dyson limits?

\textbf{Basic objects.} The paper studies Hamiltonians of the form
\[
H = H_0 + \gamma H_1,
\]
where $H_0$ is integrable, $H_1$ is a chaotic perturbation, and $\gamma$ tunes the distance from integrability.

\section*{Executive summary}
The paper gives a clean answer to a messy old problem. The authors argue that, in the crossover from integrability to chaos, the full microscopic pattern of the perturbation matrix is not what matters for level-spacing statistics. What matters is much coarser: the empirical distribution of the off-diagonal matrix elements of $H_1$ when written in the eigenbasis of the integrable part $H_0$.

From that observation they build a random-matrix surrogate
\[
\mathcal H = \mathcal D + \gamma' \mathcal M,
\]
where $\mathcal D$ is a random diagonal matrix sampled from the distribution of the eigenvalues of $H_0$, and $\mathcal M$ is a symmetric zero-diagonal random matrix whose nonzero entries are i.i.d. draws from the empirical distribution of off-diagonal entries $\langle n|H_1|m\rangle$ with $n<m$. After matching perturbation strengths through the average adjacent-gap ratio $\bar r$, this surrogate reproduces the full unfolded spacing histogram of the physical model with striking accuracy.

They validate the proposal on two very different families: perturbed spin-$1/2$ chains and quantum resonant systems. In both cases the new ensemble matches the observed crossover statistics much better than standard baselines such as the Brody fit and the Rosenzweig--Porter model. A second punchline is empirical but intriguing: the distribution of off-diagonal perturbation matrix elements often has a broad power-law regime. That suggests crossover statistics may fall into universality classes indexed by the tail shape of this matrix-element distribution rather than by a single interpolation parameter.

\section*{Problem setup and intuition}
Two benchmark limits are familiar. For integrable systems, unfolded level spacings $s$ are asymptotically Poisson,
\[
P_{\rm P}(s)=e^{-s}.
\]
For chaotic systems with time-reversal symmetry, one gets Gaussian-orthogonal-ensemble behavior, well approximated by the Wigner surmise,
\[
P_{\rm WD}(s)=\frac{\pi s}{2}e^{-\pi s^2/4}.
\]
The hard regime is the interpolation between these two limits. Existing descriptions either fit data phenomenologically or use overly rigid random-matrix models that miss system-specific intermediate shapes.

The key intuition is to write $H_1$ in the eigenbasis $\{|n\rangle\}$ of $H_0$:
\[
H_0 |n\rangle = E_n |n\rangle, \qquad H_{1,nm}=\langle n|H_1|m\rangle.
\]
The authors argue that the crossover spacing law is controlled by the statistics of the cloud of numbers $\{H_{1,nm}: n<m\}$, not by the exact geometric placement of those entries inside the matrix. Once one knows the diagonal spectrum inherited from the integrable part and the empirical distribution of off-diagonal couplings injected by the perturbation, one already has the main information needed to reproduce the transition.

\section*{Main result (informal but precise)}
The authors propose the ensemble
\[
\mathcal H = \mathcal D + \gamma'\mathcal M
\]
with the following ingredients:
\begin{itemize}
\item $\mathcal D$ is diagonal, with i.i.d. entries drawn from a smoothed version of the empirical spectrum of $H_0$.
\item $\mathcal M$ is symmetric with zero diagonal, and its off-diagonal entries are i.i.d. draws from a smoothed version of the empirical distribution of $\langle n|H_1|m\rangle$ in the $H_0$ eigenbasis.
\item $\gamma'$ is tuned so that this model has the same average adjacent-gap ratio $\bar r$ as the physical Hamiltonian $H_0+\gamma H_1$.
\end{itemize}
Then, at large Hilbert-space dimension, the unfolded nearest-neighbor spacing distribution of the surrogate matches that of the physical system across the integrability-to-chaos crossover.

The matching statistic is
\[
r_n=\frac{\min(\delta_n,\delta_{n-1})}{\max(\delta_n,\delta_{n-1})}, \qquad \delta_n = E_{n+1}-E_n,
\]
and the limiting values are $\bar r_{\rm P}=2\log 2-1\approx 0.386$ and $\bar r_{\rm WD}=4-2\sqrt 3\approx 0.536$.

\section*{Proof architecture}
\begin{enumerate}
\item \textbf{Identify the right state variables.} Start with $H=H_0+\gamma H_1$, diagonalize $H_0$, and record the samples of eigenvalues $\{E_n\}$ and off-diagonal perturbation entries $\{\langle n|H_1|m\rangle:n<m\}$.
\item \textbf{Build a surrogate ensemble.} Replace the physical Hamiltonian by a random matrix with the same coarse one-point statistics: a diagonal part sampled from the $H_0$ spectrum and an i.i.d. symmetric off-diagonal part sampled from the perturbation-entry distribution.
\item \textbf{Normalize comparisons through $\bar r$.} Tune $\gamma$ and $\gamma'$ so the physical model and surrogate sit at the same value of the average gap ratio $\bar r$.
\item \textbf{Compute unfolded spacing histograms.} Trim edge levels, unfold the remaining spectrum using a mesoscopic local-density prescription, and compare histograms for several intermediate values of $\bar r$.
\item \textbf{Benchmark against alternatives.} Compare the new ensemble to the Rosenzweig--Porter model and Brody-family fits; the proposed ensemble tracks the physical histograms substantially better.
\end{enumerate}

\section*{Why the model works}
Nearest-neighbor repulsion is a local spectral phenomenon produced by mixing of nearly resonant unperturbed states. In first approximation, that mixing is encoded by the size statistics of off-diagonal couplings, not by the exact deterministic architecture of the perturbation matrix. The diagonal spectrum of $H_0$ sets the unperturbed gaps; the distribution of $H_{1,nm}$ tells us how strongly those gaps are hybridized.

This also explains why the Rosenzweig--Porter model can miss real systems. RP hardcodes Gaussian off-diagonal statistics, whereas the physical perturbations examined here often have highly non-Gaussian, broad, approximately power-law entry distributions.

\section*{Validation, limitations, and takeaway}
For a perturbed transverse-field Ising chain with an XY-Heisenberg-type perturbation, the ensemble reproduces the physical spacing curves at several intermediate values of $\bar r$. The same strategy succeeds again for perturbed quantum resonant systems, which are structurally quite different many-body bosonic models. That cross-family agreement is what makes the proposal compelling.

A particularly interesting empirical finding is that the distribution $P_{\mathcal M}(x)$ of off-diagonal perturbation entries often shows a broad power-law regime that can be fit schematically by
\[
P_{\mathcal M}(x)\sim \frac{e^{-Cx^2}}{x^p}, \qquad 0\le p\le 1.
\]
The exact cutoff behavior seems much less important than the existence and slope of this broad regime.

The work is convincing numerically, but analytic control is still missing. The paper does not derive the crossover law from first principles or prove that all higher-order correlations become irrelevant. It also focuses on spacing statistics, not on eigenvector structure, transport, entanglement dynamics, or thermal observables.

\textbf{Bottom line.} The paper's main contribution is a new coarse-grained random-matrix principle: Poisson-to-Wigner crossover statistics are governed by the diagonal spectrum of the integrable part together with the empirical distribution of off-diagonal perturbation matrix elements.
\end{document}
